\section{Conclusion and Related Work}

This paper brings compositional principles and a typing discipline in the setting of \QBNs, addressing the lack of compositionality and modularity in the original setting with methods from denotational semantics and proof theory.
Our framework is fully compatible with classical Bayesian networks and Bayesian inference.
The crucial notion is that of \emph{quantum factors}, that directly generalize factors from the theory of (classical) Bayesian networks and satisfy the same key properties for product and sum.
\begin{center}
\begin{tabular}{|c|c|}
\hline
Classical factor over \rvs $\bX$ & Quantum factor over \rvs $\bX$ and the quantum register $\qreg$
\\
\hline
$\phi:\setofrv{\bX} \to \nnReal$	& $\someqfact: \setofrv{\bX} \to \Pos{\HH_\qreg}$
\\
\hline
\end{tabular}
\end{center}
Their development is highly non-trivial, and is our main technical contribution.
A major challenge has been to conciliate the behaviour of classical factors with quantum no-cloning.
Our product of \qfactor s shares the value of classical variables (as factors do) and coincides with tensor contraction for quantum registers, being fully compatible with both settings.

The framework enables computing the probability distribution associated to a network by multiplying its \qfactor s in any arbitrary order (not necessarily in a top-down fashion as in~\cite{GBN2014}), and by summing out irrelevant variables on sub-components.
%This feature is key to inference algorithms for (classical) Bayesian networks, looking for a most efficient order of computations.
%It is also used to study conditional independence (no-signaling) between random variables, hence it may be possible to do the same even when considering a random variable and a quantum cause, or two quantum causes.

Furthermore, we propose a typed formalism for \QBN s---based on proof-nets---that guarantees \emph{modularity} through a proof-theoretic approach.
The \emph{exact} correspondence between closed proof-nets and QBNs is guaranteed by soundness and completeness results.
Finally, it is worth noting that \qfactor s are a model of Multiplicative Linear Logic.
%Quantum Bayesian networks correspond to a particular class of proof-nets, with related semantics (\cref{th:sem_coherence}), and also with a compositional semantics (\cref{th:compositional}).
%These proof-nets are typed, and two compatible types indicate sub-graphs that can be linked together into a DAG (\cref{th:types}).
%This corresponds to linear application, which is a weak form of higher-order.
%An interesting question is whether it is possible to have a stronger form of higher-order.

\subparagraph{Related work.}

Our semantics is strongly inspired by Selinger's~\cite{Selinger2004}---the graph of the function defining a \emph{quantum factor} can be seen as an \emph{indexed tuple of matrices}, closely related to the tuples of matrices in~\cite{Selinger2004}.
Another notable denotational model built over density operators and quantum operations is~\cite{HasuoH17}.
Our typed graphical formalism based on linear logic exploits and expands a recent line of work relating Bayesian networks with \pns~\cite{EhrhardFP23,BPN} and with the information flow on type derivations~\cite{popl24}. 

As discussed in the introduction, the literature on quantum causal models (see \eg~\cite{BarrettLorenzOreshkov2019} and references therein) is vast and multifaceted.
We favor the \emph{inference} perspective propounded in~\cite{LeiferSpekkens2013}, a work which is preliminary to QBNs~\cite{GBN2014}.
While we only cite~\cite{LeiferSpekkens2013}, a substantial body of work studies quantum extension of Bayes' rule and Bayesian inference in quantum mechanics and quantum foundations.
This is however beyond the scope of our work.

Our goal is a \emph{conservative} extension of Bayesian networks to the quantum setting.
An early definition of quantum Bayesian networks---quite different from the one we consider---is given by Tucci~\cite{Tucci1995}, where amplitudes replace probabilities.
We focus on the work by Henson, Lal, and Pusey~\cite{GBN2014}, that introduces a notion of \emph{generalized Bayesian networks} including the quantum case; recent work in the same line is~\cite{KhannaAnsanelliPuseyWolfe2024}.
A closely connected formalism is that of~\cite{Fritz2012}, which focuses on (quantum) non-locality in scenarios with several quantum sources.
The results in~\cite{GBN2014} exploit also the framework of \emph{quantum networks} developed in~\cite{10.1103/PhysRevA.80.022339}.
It should be noted that \emph{quantum networks} and \emph{quantum Bayesian networks}, despite their similar names, are distinct frameworks with different purposes.
In particular, quantum networks cannot encode BNs: a classical operational-probabilistic theory in~\cite{10.1103/PhysRevA.80.022339} is not a BN---see the discussion in~\cite[Section~3.1 and Example~10]{GBN2014}.

Taking a broader perspective beyond \QBNs, we briefly mention some active research areas which share related aims. 
The development of a higher-order typed framework for quantum processes~\cite{10.1098/rspa.2018.0706,10.22331/q-2026-01-21-1978} is a timely and vibrant research line at the frontier between quantum physics and higher-order computation, exploring quantum operations on quantum operations.
Another very active line of research, concerned with compositionality and higher-order structures in both quantum and probabilistic settings, is that of string diagrams, adopting a categorical approach---see \eg~\cite{10.23638/LMCS-15(3:15)2019,10.4230/LIPIcs.MFCS.2022.80, simmons2024completelogiccausalconsistency}.
This fruitful line has a perspective different from ours.
We noted already that BNs have a dual nature, as tools for \emph{efficient probabilistic inference} and as \emph{causal models}.
While we favor the former, the line of developments based on monoidal category theory focuses on causality and is unconcerned with the computational cost of probabilistic reasoning.
As observed in~\cite[Example~9.7]{popl24}, the cost of \emph{actually} computing the semantics of a BN explodes when taking this approach, because a central notion is a product $\otimes$ which behaves like the tensor product of matrices.
We give an example in \cref{app:cost}. 

Our proof-net syntax provides an \emph{exact} characterization of QBNs, setting our work apart from other frameworks.
In the last two decades, several papers have developed variants of proof-nets accounting for quantum processes~\cite{DuncanThesis,AbramskyDuncan2009GeneralisedProofNets,esop2014,LagoFVY17,simmons2024completelogiccausalconsistency}; here we briefly discuss~\cite{simmons2024completelogiccausalconsistency} which introduces proof-nets as a tool to analyze \emph{causal} structures.
A key ingredient, which is also key in~\cite{popl24} and in our own treatment, is the (polarized) flow of information.
The framework in~\cite{simmons2024completelogiccausalconsistency} focuses on soundness and completeness (completeness being non‑trivial) \wrt\ causal string diagrams, building on Retoré’s proof-nets: causality is encoded via a non‑commutative tensor.
In contrast, we present a class of proof‑nets that is sound and complete (soundness being non‑trivial) \wrt\ QBNs.
When all data are classical, we recover BNs and~\cite{BPN}, where standard inference algorithms (\eg\ Variable Elimination, Message Passing) are available.
